\subsection{Radiometry and Surface Scattering}\label{subsec:radiometry-surface-reflectance}

This section covers the fundamental concepts of radiometry, photometry, and surface scattering functions 
that are essential for understanding light transport in rendering.

\subsubsection{Radiometry and Photometry}\label{subsubsec:radiometry-photometry}

It is important to understand the difference between radiometry (physical measurement of \gls{em} radiation)
and photometry (light as it is perceived by the human eye).
Both systems define their own set of units.
The conversion between the two is defined by the spectral luminous efficiency function $V(\lambda)$,
which peaks at approximately $555\,\mathrm{nm}$ (green).
This is why colors like green and yellow appear brighter to the human eye compared to colors like blue.

\Cref{tab:radiometric-quantities} shows the key radiometric quantities,
and \Cref{tab:photometric-quantities} lists their photometric equivalents.
Note that while rendering calculations are performed using radiometric units (specifically Radiance),
output devices (displays) and human perception operate in the photometric domain (Luminance).~\cite{tuwien_rendering_2020}.

\begin{table*}
    \centering
    \caption{Radiometric Quantities (Physical)~\cite{tuwien_rendering_2020}}
    \label{tab:radiometric-quantities}
    \begin{tabular}{lccl}
        \toprule
        \textbf{Quantity} & \textbf{Symbol} & \textbf{Unit} & \textbf{Description} \\
        \midrule
        Radiant Energy & $Q_e$ & Joule [$\mathrm{J}$] & Total energy of radiation \\
        Radiant Flux (Power) & $\Phi_e$ & Watt [$\mathrm{W}$] & Energy per unit time ($\mathrm{J/s}$) \\
        Radiant Intensity & $I_e$ & Watt/sr [$\mathrm{W/sr}$] & Flux per unit solid angle \\
        Irradiance & $E_e$ & Watt/$\mathrm{m}^2$ [$\mathrm{W/m^2}$] & Incident flux per unit area \\
        Radiance & $L_e$ & Watt/($\mathrm{sr} \cdot \mathrm{m}^2$) & Flux per unit projected area per unit solid angle \\
        \bottomrule
    \end{tabular}
\end{table*}

\begin{table*}
    \centering
    \caption{Photometric Quantities (Perceptual)~\cite{tuwien_rendering_2020}}
    \label{tab:photometric-quantities}
    \begin{tabular}{lccl}
        \toprule
        \textbf{Quantity} & \textbf{Symbol} & \textbf{Unit} & \textbf{Description} \\
        \midrule
        Luminous Energy & $Q_v$ & Talbot [$\mathrm{T}$] & Perceived energy of light \\
        Luminous Flux & $\Phi_v$ & Lumen [$\mathrm{lm}$] & Perceived power \\
        Luminous Intensity & $I_v$ & Candela [$\mathrm{cd}$] & Flux per unit solid angle ($\mathrm{lm/sr}$) \\
        Illuminance & $E_v$ & Lux [$\mathrm{lx}$] & Incident flux per unit area ($\mathrm{lm/m^2}$) \\
        Luminance & $L_v$ & Nit [$\mathrm{cd/m^2}$] & Perceived brightness of a surface \\
        \bottomrule
    \end{tabular}
\end{table*}

\subsubsection{Bidirectional Reflectance Distribution Function (BRDF)}\label{subsubsec:brdf}

The \gls{brdf} defines how a material scatters incoming light from direction $\omega_i$ 
into outgoing direction $\omega_o$ at a given surface point $x$~\cite{136913.136929}.
Formally the \gls{brdf} is the differential ratio of reflected radiance $\mathrm{d}L_r$ 
to incident irradiance $\mathrm{d}E_i$:

\begin{align}
    f_r(x, \omega_i, \omega_o) &= \frac{\mathrm{d} L_r(x, \omega_o)}{\mathrm{d} E_i(\omega_i)} \\
    &= \frac{\mathrm{d} L_r(x, \omega_o)}{L_o(y_i, -\omega_i) \, |\cos\theta_i| \, \mathrm{d}\omega_i}
\end{align}

The unit of the \gls{brdf} is radiance divided by irradiance, which is $\mathrm{sr^{-1}}$ (inverse steradians).

In terms of fundamental flux quantities, the \gls{brdf} can also be expressed as:

\begin{align}
    f_r(x, \omega_i, \omega_o) = \frac{\mathrm{d}^3 \Phi_r}{\mathrm{d}^2 \Phi_i \, \mathrm{d}\omega_o \, |\cos\theta_o|}
\end{align}

The \gls{brdf} must also satisfy the following requirements:
\begin{itemize}
    \item \textbf{Positivity}: Light cannot be negative, so it must satisfy $f_r(x, \omega_i, \omega_o) \geq 0$ for all $\omega_i, \omega_o$
    \item \textbf{Helmholtz Reciprocity}: States that a ray and its reverse ray must follow the same optical path 
        (e.g.\ reflections, refractions) $f_r(x, \omega_i, \omega_o) = f_r(x, \omega_o, \omega_i)$
    \item \textbf{Energy Conservation}: The cosine weighted integral of the \gls{brdf} over the hemisphere must be
            1 or less for all $\omega_o$:
    \[
        \int_{\Omega} f_r(x, \omega_i, \omega_o) |\cos\theta_o| \mathrm{d}\omega_o \leq 1
    \]
\end{itemize}

\subsubsection{Bidirectional Transmittance Distribution Function (BTDF)}\label{subsubsec:btdf}

While the \gls{brdf} covers reflection, the \gls{btdf} handles light that passes through a surface (transmission).
The lower hemisphere $\Omega^-$ is used for incoming light (\gls{brdf}), 
and the upper hemisphere $\Omega^+$ is used for outgoing light (\gls{btdf}).
Unlike reflection, the \gls{btdf} is not symmetric, as light can be transmitted in one direction but not the other,
and when light transmits across an interface between two media, with different refractive indices $n_1$ and $n_2$,
the \gls{btdf} is scaled by the ratio of the squared indices of refraction ${({n_2}/{n_1})}^2$ 
according to Snell's law.~\cite{136913.136929,Born_Wolf_Bhatia_Clemmow_Gabor_Stokes_Taylor_Wayman_Wilcock_1999}

\[
    f_t(\omega_i, \omega_t) = \frac{d L_t(\omega_t)}{d E_i(\omega_i)} = \frac{d L_t(\omega_t)}{L_i(\omega_i) |\cos\theta_i| d\omega_i}
\]


\subsubsection{Bidirectional Scattering Distribution Function (BSDF)}\label{subsubsec:bsdf}

The \gls{bsdf} unifies reflection and transmission into a single function
over the entire sphere of directions $4\pi$ steradians,
logically representing the sum of $f_r$ and $f_t$.~\cite{Bartell1980BSDF,Pharr2016PBR}

\[
    f_s(\omega_i, \omega_o) = f_r(\omega_i, \omega_o) + f_t(\omega_i, \omega_t)
\]

Depending on whether the incident and outgoing directions are on the same side 
of the surface (reflection) or opposite sides (transmission), 
the BSDF decomposes into two distinct components,
because light cannot simultaneously reflect and transmit 
along the exact same outgoing vector $\omega_o$ at a sharp boundary point.
Therefore, the general BSDF is defined piecewise as:

\begin{align}
    f_s(x, \omega_i, \omega_o) = 
    \begin{cases} 
      f_r(x, \omega_i, \omega_o), & \text{reflection} \\ 
      f_t(x, \omega_i, \omega_o), & \text{transmission}
    \end{cases}
\end{align}

\subsubsection{Isotropic vs Anisotropic Scattering}\label{subsubsec:isotropic-anisotropic-scattering}

Scattering at a point depends on 4 angular dimensions: 
the incident direction $(\theta_i, \phi_i)$ and the outgoing direction $(\theta_o, \phi_o)$.
~\cite{Pharr2016PBR,10.1145/133994.134078}

\paragraph{Isotropic Scattering:}

A surface is isotropic if its scattering properties are invariant under rotation about the surface normal $n$.
Spinning the object around its normal does not change the appearance of the surface.
Therefore the scattering function no longer depends on the azimuthal angles $\phi_i$ and $\phi_o$,
and the \gls{bsdf} reduces to a function of only the polar angles $\theta_i$ and $\theta_o$: 
$f_s(\theta_i, \theta_o, \Delta\phi)$
Examples of isotropic materials include Lambertian surfaces, metals, dielectrics, glass, and some plastics.

\paragraph{Anisotropic Scattering:}

Anisotropic surfaces have scattering properties that vary with rotation about the surface normal.
This leaves the \gls{bsdf} as a 4D function of both the polar and azimuthal angles:
$f_s(\theta_i,\phi_i,\theta_o,\phi_o)$.
Anisotropy commonly results from an oriented microscopic structure.
For example mie scattering from aligned fibers, water droplets, brushed metals, or dust.

In practice, anisotropic models often use a tangent direction $t$ in addition to the surface normal
$n$ to describe this orientation. The tangent and bitangent define a local frame in which the
microfacets or scattering lobes can have different roughness values along different axes. When the
roughness is equal in both tangent directions, the model reduces to the isotropic case.




